\section{结论与建议}

\subsection{三项任务的主要结论}

三项任务的主要结论、支撑证据与适用边界汇总于表~\ref{tab:conclusions}。
每条结论都指向支撑它的图表编号，并注明适用边界。

\begin{table}[H]
\centering\small
\caption{主要结论、支撑证据与适用边界}
\label{tab:conclusions}
\begin{tabular}{p{4.2cm}p{3.4cm}p{4.6cm}}
\toprule
结论 & 支撑证据 & 适用边界 \\
\midrule
化学特征对品评小组均分有可复现的预测能力，主模型等级 MAE 0.479，显著优于只知风格的基线 0.527 & 表~\ref{tab:q1-main}、图~\ref{fig:q1-dim} & 连续 $R^2$ 只有 0.184，信度中位数 0.53，预测能力有限 \\
\addlinespace
只用附件无法还原原始等级，两阶段对照 C1 等级 MAE 0.786 & 表~\ref{tab:q1-main}、表~\ref{tab:prep-reference} & C1 按自身规则选定等级映射，仍以同一原始等级为真值 \\
\addlinespace
多数维度可靠支持的关键化合物是 1 到 2 个，第 3 位以报告为主 & 表~\ref{tab:q2-rep}、图~\ref{fig:q2-stability} & Top-3 是关联结果，不是因果主张 \\
\addlinespace
选择规则由模拟还原实验确定，只用化学矩阵，不依赖感官数据 & 图~\ref{fig:q2-recovery} & 模拟假定少数化合物按秩线性可加，与真实结构存在差异 \\
\addlinespace
风格条件生成模型在留出啤酒上通过大部分验收项，缺陷约束在几乎所有风格对上起作用 & 表~\ref{tab:q3-gen}、图~\ref{fig:q3-fingerprint} & 验收阈值在预运行后固定，属工程标准，不是独立验证 \\
\addlinespace
迁移只沿一条线段搜索，$\alpha^*$ 小不说明不存在满足约束的目标风格谱 & 图~\ref{fig:q3-transfer-fig} & 结论针对选定路径与问题一模型 \\
\addlinespace
问题一主输出向中间等级收缩，几乎不给出 5 到 6 级，第二输出保持等级分布 & 表~\ref{tab:q1-rules}、图~\ref{fig:q1-grade} & 这是以等级 MAE 为主指标的直接后果 \\
\bottomrule
\end{tabular}
\end{table}

三个问题共享同一批啤酒与化学特征，问题一的折外预测与问题二的共表达模块分别作为问题三感官画像与生成模型的输入，形成完整的方法链条。
问题一与问题二的结论以验收过的量化结果为支撑，问题三的结论待正式结果验收后同步更新。

\subsection{面向风味分析与酿造试验的建议}

化学数据对品评小组均分的解释力有限，连续 $R^2$ 为 0.184，信度中位数为 0.53，5 到 6 级几乎无法由化学谱预测。
因此风味分析不应把 0 到 6 级感官评分当作化学谱的确定性函数，高分啤酒的判定仍需品评小组，化学预测适用于筛查与排序。

正文的关键化合物是关联结果，不是因果主张。
酿造试验若要以关键化合物作为配方依据，宜对单一化合物做剂量与响应方向的加标，逐项确认作用方向与有效范围。

题面把缺陷定义为阈值以上，阈值以下的化合物不构成缺陷，迁移约束因此只在超阈值项上生效。
低于检出限的浓度只能给出判断范围，不能作为精确靶值。
生成样本围绕风格均值分布，适合作为配方探索的起点，不能替代真实酿造与品评。
